\section{Iteratively Computing Approximate Lewis Weights}
\label{sec:compute}

Pseudocode of our algorithm is given in Figure~\ref{fig:iteration}.

\begin{figure}[ht]

\begin{algbox}

$\ww = \textsc{LewisIterate}(\AA, p, \beta, \ww)$

\begin{enumerate}
	\item For $i = 1 \ldots n$
	\begin{enumerate}
		\item Let $\tilde{\ttau_i} \approx_{\beta} \ttau_{i} \left( \WW^{1/2 - 1/p} \AA \right)$ be a $\beta$-approximation of the statistical leverage score of row $i$ in $\WW^{1/2 - 1/p} \AA$.
		\item Set $\wwhat_i \leftarrow \left( \ww_i^{2/p - 1} \tilde{\ttau}_i \right)^{p/2} \approx_{\beta^{p/2}} (\aa_i^T \left( \AA^T \WW^{1 - 2/p} \AA \right)^{-1} \aa_i)^{p/2}$.
	\end{enumerate}
	\item Return $\wwhat$.
\end{enumerate}

$\ww =\textsc{ApproxLewisWeights}(\AA, p, \beta, T)$

\begin{enumerate}
	\item Initialize $\ww_i  = 1$
	\item For $t = 1 \ldots T$
	\begin{enumerate}
		\item Set $\ww \leftarrow \textsc{LewisIterate}(\AA, p, \beta, \ww)$.
	\end{enumerate}
	\item Return $\ww$.
\end{enumerate}
\end{algbox}

\caption{Iterative Algorithm for Computing Lewis Weights}
\label{fig:iteration}

\end{figure}

In our proof, we make use of the generalization of approximations
 to the matrix setting via the Loewner partial ordering.
For two symmetric matrices $\PP$ and $\QQ$, we have $\PP \preceq \QQ$
if $\PP - \QQ$  is positive semi-definite.
We will then use $\PP \approx_{\alpha} \QQ$ to denote
\begin{equation*}
\frac{1}{\alpha} \PP \preceq \QQ \preceq \alpha \PP. 
\end{equation*}
Note that this is equivalent to $\xx^T \PP \xx \approx_{\alpha} \xx^T \QQ \xx$ for
all vectors $\xx$.
Two facts of particular importance for spectral approximation of matrices
are its composition and inversion.
\begin{fact}
\label{fact:operatorApprox}
\begin{enumerate}
\item If $\PP \approx_{\alpha} \QQ$, then for any matrix $\AA$, we also have
$\AA^T \PP \AA \approx_{\alpha} \AA^T \QQ \AA$.
\item If $\PP \approx_{\alpha} \QQ$, then $\PP^{-1} \approx_{\alpha} \QQ^{-1}$.
\end{enumerate}
\end{fact}

First, we show that (for $p < 4$) $\textsc{LewisIterate}(\AA, p, 1, _)$ acts as a
contraction mapping with respect to the $\ell_\infty$ distance of the $\log$ weights:
\begin{lemma}
\label{lem:contraction}
Given $\AA$, $p$, and weight sets $\vv$ and $\ww$ such that $\vv \approx_{\alpha} \ww$,
\begin{equation*}
\textsc{LewisIterate}(\AA, p, 1, \vv) \approx_{\alpha^{|p/2 - 1|}} \textsc{LewisIterate}(\AA, p, 1, \ww).
\end{equation*}
\end{lemma}
\begin{proof}
Since $\vv \approx_{\alpha} \ww$, we have, in the matrix setting
\begin{equation*}
\VV^{1 - 2/p} \approx_{\alpha^{|1 - 2/p|}} \WW^{1 - 2/p}.
\end{equation*}
Applying both items of Fact~\ref{fact:operatorApprox} we then have
\begin{equation*}
\left( \AA^T \VV^{1 - 2/p} \AA \right)^{-1} \approx_{\alpha^{|1-2/p|}} \left( \AA^T \WW^{1 - 2/p} \AA \right)^{-1}.
\end{equation*}
Then applying the definition of operator approximation, we have, for all $i$,
\begin{equation*}
\aa_i^T \left( \AA^T \VV^{1 - 2/p} \AA \right)^{-1} \aa_i \approx_{\alpha^{|1-2/p|}} \aa_i^T \left( \AA^T \WW^{1 - 2/p} \AA \right)^{-1} \aa_i.
\end{equation*}
Taking the $p/2$ power of both sides gives the desired result.
\end{proof}
An immediate corollary of this lemma is
\begin{corollary}
\label{cor:convergence}
Given $\AA$ with $\ell_p$ Lewis weights $\wwbar$ and a set of weights $\ww$ with $\ww \approx_{\alpha} \wwbar$,
\begin{equation*}
\textsc{LewisIterate}(\AA, p, \beta, \ww) \approx_{\beta^{p/2} \alpha^{|p/2 - 1|}} \wwbar.
\end{equation*}
\end{corollary}
\begin{proof}
This follows from applying Lemma~\ref{lem:contraction} to $\ww$ and $\wwbar$, noting that
$\textsc{LewisIterate}(\AA, p, \beta, \ww) \approx_{\beta^{p/2}} \textsc{LewisIterate}(\AA, p, 1, \ww)$ and that
$\textsc{LewisIterate}(\AA, p, 1, \wwbar) = \wwbar$.
\end{proof}

Interestingly, another corollary is that for $p < 4$, the Lewis weights exist and are unique:
\begin{corollary}
\label{cor:itlewis}
For all $p < 4$, there exists a unique assignment of weights $\ww$ that are a fixed point of
$\textsc{LewisIterate}(\AA, p, 1, _)$, or equivalently that satisfy Definition~\ref{def:lewis}.
\end{corollary}
\begin{proof}
This is just the Banach fixed point theorem applied to $\textsc{LewisIterate}(\AA, p, 1, _)$:
whenever $p < 4$, $|1-2/p| < 1$, so the iteration is a contraction mapping and has a unique
fixed point.
\end{proof}

We can also show that after one step $\ww$ is already polynomially close to $\wwbar$:
\begin{lemma}
\label{lem:initialization}
After $t = 1$ in $\textsc{ApproxLewisWeights}(\AA, p, \beta, T)$, $\ww_i \approx_{\beta^{p/2} n^{|p/2 - 1|}} \wwbar_i$.
\end{lemma}
\begin{proof}
We first note that we may assume, without loss of generality, that $\AA^T \WWbar^{1 - 2/p} \AA = \II$.
This is because both the definition of Lewis weights and the algorithms are invariant under replacing
$\AA$ with $\AA \RR$, with $\RR$ any full-rank $d$ by $d$ matrix.

Then we claim that
\begin{equation*}
\AA^T \AA \approx_{n^{|1 - 2/p|}} I.
\end{equation*}

First, we let $\bb_i = \wwbar_i^{1/2 - 1/p} \aa_i$, so that $\aa_i = \wwbar_i^{1/p - 1/2} \bb_i$.
We note that by the definition of leverage score and Lewis weight, $\norm{\bb_i}_2 = \ww_i$.

Now, consider any unit vector $\uu$.  We have
\begin{align*}
1 &= \uu^T \uu \\
&= \uu^T \BB^T \BB \uu \\
&= \sum_i (\uu_i^T \bb_i)^2 \\
&= \sum_i \wwbar_i (\wwbar_i^{-1} (\uu_i^T \bb_i)^2).
\end{align*}

On the other hand, we have
\begin{align*}
\uu^T \AA^T \AA \uu &= \sum_i \wwbar_i^{2/p - 1} (\uu_i^T \bb_i)^2 \\
&= \sum_i \wwbar_i^{2/p} (\wwbar_i^{-1} (\uu_i^T \bb_i)^2).
\end{align*}

Furthermore, $\sum_i \wwbar_i^{-1} (\uu_i^T \bb_i)^2 \leq n$, since each term is
$\leq 1$ (as, by Cauchy-Schwarz, $(\uu_i^T \bb_i)^2 \leq \norm{\bb_i}_2^2 = \wwbar_i$).

Then the worst-case distortion would be if $\sum_i \wwbar_i^{-1} (\uu_i^T \bb_i)^2 = n$ and all $\wwbar_i$ are $\frac{1}{n}$.  In
that case, the distortion is $n^{|1 - 2/p|}$, as desired.

Finally, the $\ww$ after one step are
\begin{equation*}
\ww_i \approx_{\beta^{p/2}} (\aa_i^T (\AA^T \AA)^{-1} \aa_i)^{p/2}.
\end{equation*}
These are then at most $\beta^{p/2} n^{p/2 |1-2/p|} = \beta^{p/2} n^{|p/2 - 1|}$ off from $\wwbar_i$.
\end{proof}

Combining this initial condition with the convergence result allows us to bound
the total number of steps, giving a proof of Theorem~\ref{thm:main}.

\begin{proof}[Proof of Lemma~\ref{lem:computeLewis}]
The total multiplicative contribution of the blowups from $\beta$ is at most $\beta^{\frac{p/2}{1 - |p/2 - 1|}}$.
The contribution from the starting error is at most $n^{|p/2 - 1|^T}$.

Then if $\beta \leq n^{\theta \frac{1 - |p/2 - 1|}{p}}$ and $T \geq \frac{\log(2 / \theta)}{1 - |p/2 - 1|}$,
each provides at most $n^{\theta / 2}$ error, so the result is an $n^{\theta}$-approximation.
\end{proof}

